% =====================================================================
\section{Máquinas térmicas}
% =====================================================================
La primera ley no prohíbe que el calor pase de un cuerpo frío a uno caliente ni que un motor convierta
íntegro en trabajo el calor de un único foco; la experiencia dice que no ocurre. La segunda ley fija el
\emph{sentido} de los procesos y \emph{limita} el trabajo útil \cite{gonzalez1,gonzalez2}. Su terreno natural
son las máquinas cíclicas: en un ciclo $\Delta U=0$, luego $Q_{\mathrm{ciclo}}=-W_{\mathrm{ciclo}}$.

\begin{definicion}[Motor, refrigerador y bomba de calor]
Sean dos focos a $T_h>T_c$ y una sustancia que recorre un ciclo; los signos son los del sistema
(figura~\ref{fig:maquinas}).
\begin{itemize}[leftmargin=1.4em]
\item Un \emph{motor} absorbe $Q_h>0$ del foco caliente, cede $Q_c<0$ al frío y entrega trabajo, $W<0$.
\item Un \emph{refrigerador} extrae $Q_c>0$ del foco frío, recibe trabajo $W>0$ y cede $Q_h<0$ al caliente.
\item Una \emph{bomba de calor} es el mismo dispositivo que el refrigerador, pero su objetivo es el calor
  $-Q_h$ que entrega al foco caliente.
\end{itemize}
\end{definicion}

\begin{definicion}[Rendimiento y coeficientes de operación]
El rendimiento es la energía útil entre la energía que hay que aportar. Con $\Delta U=0$, $W=-(Q_h+Q_c)$ y
\begin{equation}\label{eq:rendimientos}
  \eta\equiv\frac{-W}{Q_h}=1+\frac{Q_c}{Q_h},\qquad
  \varepsilon\equiv\frac{Q_c}{W}=\frac{1}{-1-Q_h/Q_c},\qquad
  \nu\equiv\frac{-Q_h}{W}=\varepsilon+1 .
\end{equation}
Si el ciclo tiene más de dos tramos con calor, $Q_h$ se sustituye por el \emph{calor total absorbido}
$Q_{\mathrm{abs}}=\sum_{Q_i>0}Q_i$, y $\eta=|W_{\mathrm{neto}}|/Q_{\mathrm{abs}}$.
\end{definicion}

\begin{figure}[htbp]
\centering
\begin{tikzpicture}[font=\small,
   foco/.style={draw=navy,fill=fondomedio,minimum width=2.7cm,minimum height=0.62cm,inner sep=2pt},
   maq/.style={draw=navy,fill=white,circle,minimum size=0.95cm,inner sep=0pt,font=\small\bfseries},
   fl/.style={-{Stealth[length=5pt]},line width=1pt}]
  \foreach \x/\nom/\tit/\dir in {0/M/Motor/1,4.6/R/Refrigerador/-1,9.2/B/Bomba de calor/-1}{
    \node[foco] (h\nom) at (\x,2.4) {foco caliente $T_h$};
    \node[foco] (c\nom) at (\x,-0.4) {foco frío $T_c$};
    \node[maq] (m\nom) at (\x,1.0) {\nom};
    \node[font=\small\bfseries,navy] at (\x,3.25) {\tit};
  }
  % motor
  \draw[fl,rojo] (hM.south) -- (mM.north) node[midway,left=1pt] {$Q_h>0$};
  \draw[fl,acento] (mM.south) -- (cM.north) node[midway,left=1pt] {$Q_c<0$};
  \draw[fl,navy] (mM.east) -- ++(1.0,0) node[right] {$W<0$};
  % refrigerador
  \draw[fl,rojo] (mR.north) -- (hR.south) node[midway,left=1pt] {$Q_h<0$};
  \draw[fl,acento] (cR.north) -- (mR.south) node[midway,left=1pt] {$Q_c>0$};
  \draw[fl,navy] ($(mR.east)+(1.0,0)$) node[right] {$W>0$} -- (mR.east);
  % bomba
  \draw[fl,rojo,line width=1.5pt] (mB.north) -- (hB.south) node[midway,left=1pt] {$Q_h<0$};
  \draw[fl,acento] (cB.north) -- (mB.south) node[midway,left=1pt] {$Q_c>0$};
  \draw[fl,navy] ($(mB.east)+(1.0,0)$) node[right] {$W>0$} -- (mB.east);
  \node[font=\footnotesize,gris] at (0,-1.1) {objetivo: $-W$};
  \node[font=\footnotesize,gris] at (4.6,-1.1) {objetivo: $Q_c$};
  \node[font=\footnotesize,gris] at (9.2,-1.1) {objetivo: $-Q_h$};
\end{tikzpicture}
\caption{Las tres máquinas cíclicas entre dos focos, con los signos del sistema. Motor: $\eta=-W/Q_h$;
refrigerador: $\varepsilon=Q_c/W$; bomba de calor: $\nu=-Q_h/W$.}
\label{fig:maquinas}
\end{figure}

\begin{observacion}[Un COP mayor que 1 no viola nada]
Un refrigerador que consume \SI{100}{J} y extrae \SI{300}{J} del foco frío tiene $\varepsilon=3$ y cede
$-Q_h=\SI{400}{J}$, de modo que como bomba $\nu=4$. El trabajo no se convierte en calor extraído: solo lo
\emph{transporta}.
\end{observacion}

% =====================================================================
\section{Segunda ley y entropía}
% =====================================================================

\subsection{Enunciados de la segunda ley}

\begin{teorema}[Enunciados equivalentes de la segunda ley]
\begin{enumerate}[leftmargin=1.6em,label=(\arabic*)]
\item \emph{Clausius.} Es imposible un dispositivo cíclico cuyo único efecto sea pasar calor de un foco a otro
  de mayor temperatura.
\item \emph{Kelvin--Planck.} Es imposible un dispositivo cíclico cuyo único efecto sea extraer calor de un foco
  y convertirlo íntegro en trabajo (móvil perpetuo de segunda especie).
\item \emph{Carnot.} Ningún motor que opere entre dos focos tiene un rendimiento mayor que el del ciclo de
  Carnot entre ellos.
\item \emph{Aumento de la entropía.} La entropía de un sistema aislado nunca disminuye: $\Delta S\geq0$.
\end{enumerate}
\end{teorema}

La equivalencia se demuestra en cadena, $(1)\Rightarrow(4)\Rightarrow(3)\Rightarrow(2)\Rightarrow(1)$, negando
cada enunciado y construyendo una máquina que viola el anterior \cite{gonzalez1}. El enunciado (4) es el
general: no se limita a ciclos y es el que se usa en los problemas.

\subsection{Entropía}

\begin{definicion}[Entropía]
El calor reversible es una diferencial inexacta con factor integrante $1/T$: $\dbar Q_{\mathrm{rev}}/T$ es
exacta y define la \emph{entropía} $S$, función de estado extensiva en \si{J/K}:
\begin{equation}\label{eq:dS}
  \dd S=\frac{\dbar Q_{\mathrm{rev}}}{T},\qquad
  \Delta S=S_B-S_A=\int_{A,\ \mathrm{rev}}^{B}\frac{\dbar Q}{T}.
\end{equation}
\end{definicion}

\begin{proposicion}[Exactitud de $\dbar Q_{\mathrm{rev}}/T$]
Para un sistema hidrostático, $\dbar Q_{\mathrm{rev}}/T$ es exacta si y solo si
\begin{equation}\label{eq:energetica}
  \pd UVT=T\pd pTV-p,
\end{equation}
la \emph{ecuación termodinámica de la energía}. Para el gas ideal, $T(\partial p/\partial T)_V=nRT/V=p$ y
$\pdi UVT=0$: la ley de Joule es una consecuencia de la segunda ley.
\end{proposicion}

\begin{demostracion}
Con $\dbar Q_{\mathrm{rev}}=\dd U+p\,\dd V$ y $\dd U=C_V\dd T+\pdi UVT\dd V$,
\[
  \frac{\dbar Q_{\mathrm{rev}}}{T}=\underbrace{\frac{C_V}{T}}_{M}\dd T+\underbrace{\frac1T\Bigl[\pd UVT+p\Bigr]}_{N}\dd V .
\]
Derivadas cruzadas: $\partial_VM=\frac1T\frac{\partial^2U}{\partial V\partial T}$ y
$\partial_TN=\frac1T\Bigl[\frac{\partial^2U}{\partial T\partial V}+\pd pTV\Bigr]-\frac{1}{T^2}\Bigl[\pd UVT+p\Bigr]$.
Las segundas derivadas de $U$ son iguales (Schwarz) y se cancelan; igualar $\partial_VM=\partial_TN$ deja
$\frac1T\pdi pTV=\frac1{T^2}\bigl[\pdi UVT+p\bigr]$, que es \eqref{eq:energetica}.
\end{demostracion}

\begin{teorema}[Desigualdad de Clausius y balance de entropía]
En un ciclo cualquiera, con $T_F$ la temperatura del foco en cada intercambio, $\oint\dbar Q/T_F\leq0$, con
igualdad solo si el ciclo es reversible. Para un proceso cualquiera entre $A$ y $B$,
\begin{equation}\label{eq:balance}
  \Delta S=\int_A^B\frac{\dbar Q}{T_F}+S_{\mathrm{gen}},\qquad S_{\mathrm{gen}}\geq0,
\end{equation}
donde el primer término es la entropía que transporta el calor y $S_{\mathrm{gen}}$ la que producen las
irreversibilidades (fricción, disipación, expansión libre, diferencias finitas de $T$).
\end{teorema}

\begin{proposicion}[Principio de aumento de la entropía]
Para el universo termodinámico (sistema más entorno, aislado en conjunto),
\begin{equation}\label{eq:universo}
  \Delta S_{\mathrm{univ}}=\Delta S_{\mathrm{sis}}+\Delta S_{\mathrm{ent}}=S_{\mathrm{gen}}\geq0,
\end{equation}
con igualdad si y solo si todo el proceso es reversible. En un proceso adiabático $\Delta S_{\mathrm{sis}}=S_{\mathrm{gen}}$:
reversible, $\Delta S=0$ (isoentrópico); irreversible, $\Delta S>0$.
\end{proposicion}

\begin{observacion}[$Q=0$ no implica $\Delta S=0$]
En la expansión libre de un gas ideal ($V\to2V$, recipiente rígido y adiabático) $Q=W=\Delta U=0$ y
$T_2=T_1$, pero $\Delta S=nR\ln2>0$ ($\SI{5,76}{J/K}$ para un mol): toda esa entropía es generada.
\end{observacion}

% =====================================================================
\section{Cálculo de cambios de entropía}
% =====================================================================

\begin{recuadro}{Regla de oro}
$S$ es función de estado: $\Delta S_{\mathrm{sis}}$ solo depende de los estados inicial y final. Si el proceso
real es irreversible, se calcula $\Delta S$ por \emph{cualquier camino reversible} que una los mismos estados.
El calor del proceso real solo sirve para la entropía del \emph{entorno}.
\end{recuadro}

\subsection{Sistema: gas ideal y cuerpos}

\begin{proposicion}[Entropía del gas ideal]
\begin{equation}\label{eq:Sgas}
  \Delta S=\int_{T_1}^{T_2}\frac{C_V}{T}\,\dd T+nR\ln\frac{V_2}{V_1}
          =\int_{T_1}^{T_2}\frac{C_p}{T}\,\dd T-nR\ln\frac{p_2}{p_1},
\end{equation}
y para el gas perfecto las integrales valen $C_V\ln(T_2/T_1)$ y $C_p\ln(T_2/T_1)$.
\end{proposicion}

\begin{demostracion}
Con $\dd U=C_V\dd T$ y $p=nRT/V$: $\dd S=(\dd U+p\,\dd V)/T=C_V\frac{\dd T}{T}+nR\frac{\dd V}{V}$. Para la segunda
forma se sustituye $\frac{\dd V}{V}=\frac{\dd T}{T}-\frac{\dd p}{p}$ (de $\ln V=\ln nR+\ln T-\ln p$) y se usa
$C_V+nR=C_p$.
\end{demostracion}

\begin{table}[htbp]
\centering
\caption{Cambios de entropía que aparecen en los seminarios.}
\label{tab:dS}
\small
\renewcommand{\arraystretch}{1.75}
\begin{tabular}{ll}
\toprule
Sistema o proceso & $\Delta S$\\
\midrule
Gas perfecto, isotermo (o expansión libre) & $nR\ln\dfrac{V_2}{V_1}=nR\ln\dfrac{p_1}{p_2}$\\
Gas perfecto, isócoro / isóbaro & $C_V\ln\dfrac{T_2}{T_1}$\ /\ $C_p\ln\dfrac{T_2}{T_1}$\\
Gas perfecto, adiabático reversible & $0$\\
Cuerpo de capacidad $C$ constante, $T_1\to T_2$ & $C\ln\dfrac{T_2}{T_1}$\\
Foco térmico a $T_F$ que recibe $Q_F=-Q_{\mathrm{sis}}$ & $\dfrac{Q_F}{T_F}=-\dfrac{Q_{\mathrm{sis}}}{T_F}$\\
Foco mecánico, fuente de trabajo eléctrico & $0$\\
\bottomrule
\end{tabular}
\end{table}

\subsection{Entorno: focos y cuerpos finitos}

Un \emph{foco térmico} no cambia de temperatura, así que su calor es reversible desde su lado:
$\Delta S_F=Q_F/T_F$, con $Q_F=-Q_{\mathrm{sis}}$. Un \emph{cuerpo finito} sí cambia de temperatura, y su entropía se
calcula como la de un sistema, $C\ln(T_f/T_i)$. Las fuentes de trabajo (foco mecánico, generador) no cambian
de entropía; una resistencia o un freno que disipa trabajo sí la produce.

\begin{ejemplo}[Tres casos tipo]
\emph{(i) Calor entre focos}: \SI{1,20}{kJ} pasan de un foco a \SI{400}{K} a otro a \SI{300}{K}:
$\Delta S_{\mathrm{univ}}=-\frac{1200}{400}+\frac{1200}{300}=-3+4=\SI{1}{J/K}>0$.
\emph{(ii) Dos cuerpos iguales} ($C=\SI{1}{kJ/K}$) a 400 y \SI{300}{K} en contacto: $T_f=\SI{350}{K}$ y
$\Delta S_{\mathrm{univ}}=C\ln\frac{350^2}{400\cdot300}=C\ln\frac{49}{48}=\SI{20,6}{J/K}$.
\emph{(iii) Disipación}: una resistencia que entrega $W$ a un cuerpo rígido y adiabático de capacidad $C$
produce $\Delta S_{\mathrm{univ}}=C\ln(T_f/T_i)$ con $C(T_f-T_i)=W$.
\end{ejemplo}

% =====================================================================
\section{Ciclo y teorema de Carnot}
% =====================================================================

\begin{definicion}[Ciclo de Carnot]
Dos isotermas reversibles a $T_h$ y $T_c$ unidas por dos adiabáticas reversibles: $1\to2$ isoterma a $T_h$
(absorbe $Q_h$), $2\to3$ adiabática, $3\to4$ isoterma a $T_c$ (cede $Q_c$), $4\to1$ adiabática
(figura~\ref{fig:carnot}).
\end{definicion}

\begin{proposicion}[Rendimiento de Carnot]
Para cualquier sustancia,
\begin{equation}\label{eq:carnot}
  \eta_C=1-\frac{T_c}{T_h},\qquad \varepsilon_C=\frac{T_c}{T_h-T_c},\qquad \nu_C=\frac{T_h}{T_h-T_c}.
\end{equation}
\end{proposicion}

\begin{demostracion}
En el plano $T$--$S$ el ciclo es un rectángulo: las adiabáticas reversibles son isoentrópicas
($\Delta S=0$, verticales). Con $\Delta S=S_2-S_1=S_3-S_4$,
$Q_h=\int T\,\dd S=T_h\Delta S$ y $Q_c=-T_c\Delta S$, luego $\eta_C=1+Q_c/Q_h=1-T_c/T_h$. No se ha usado
ninguna ecuación de estado. Para un gas ideal en el plano $p$--$V$ se llega a lo mismo: las adiabáticas
dan $T_hV_2^{\gamma-1}=T_cV_3^{\gamma-1}$ y $T_hV_1^{\gamma-1}=T_cV_4^{\gamma-1}$, de donde $V_2/V_1=V_3/V_4$ y
$Q_c/Q_h=\frac{nRT_c\ln(V_4/V_3)}{nRT_h\ln(V_2/V_1)}=-\frac{T_c}{T_h}$.
\end{demostracion}

\begin{figure}[htbp]
\centering
\begin{subfigure}[b]{0.49\textwidth}\centering
\begin{tikzpicture}
% T_h = 1,5 ; T_c = 1 ; gamma = 5/3 ; V1 = 1, V2 = 2,5, V3 = 2,5*1,5^1,5 = 4,593, V4 = 1,5^1,5 = 1,837
\begin{axis}[graf,width=6.9cm,height=5.4cm,xmin=0.5,xmax=5.2,ymin=0,ymax=1.75,
    xtick=\empty,ytick=\empty,xlabel={$V$},ylabel={$p$},clip=false]
  \fill[fondomedio] plot[domain=1:2.5,samples=50] (\x,{1.5/\x})
    -- plot[domain=2.5:4.593,samples=60] (\x,{0.6*(2.5/\x)^(5/3)})
    -- plot[domain=4.593:1.837,samples=60] (\x,{1/\x})
    -- plot[domain=1.837:1,samples=60] (\x,{0.54433*(1.837/\x)^(5/3)}) -- cycle;
  \addplot[rojo,line width=1.1pt,domain=1:2.5,samples=60] {1.5/x};
  \addplot[navy,line width=1.1pt,domain=2.5:4.593,samples=60] {0.6*(2.5/x)^(5/3)};
  \addplot[acento,line width=1.1pt,domain=1.837:4.593,samples=60] {1/x};
  \addplot[navy,line width=1.1pt,domain=1:1.837,samples=60] {0.54433*(1.837/x)^(5/3)};
  \draw[rojo,-{Stealth[length=5pt]}] (axis cs:1.5,1.0) -- (axis cs:1.75,{1.5/1.75});
  \draw[acento,-{Stealth[length=5pt]}] (axis cs:3.6,{1/3.6}) -- (axis cs:3.2,{1/3.2});
  \fill[black] (axis cs:1,1.5) circle (1.5pt) node[anchor=south east,font=\scriptsize,inner sep=2pt] {1};
  \fill[black] (axis cs:2.5,0.6) circle (1.5pt) node[anchor=south west,font=\scriptsize,inner sep=2pt] {2};
  \fill[black] (axis cs:4.593,0.21772) circle (1.5pt) node[anchor=south west,font=\scriptsize,inner sep=2pt] {3};
  \fill[black] (axis cs:1.837,0.54433) circle (1.5pt) node[anchor=north east,font=\scriptsize,inner sep=2pt] {4};
  \node[font=\scriptsize,rojo,anchor=south west] at (axis cs:1.75,0.92) {$T_h$, $Q_h>0$};
  \node[font=\scriptsize,acento,anchor=north] at (axis cs:3.0,0.27) {$T_c$, $Q_c<0$};
  \node[font=\scriptsize,navy] at (axis cs:2.3,0.54) {$|W|$};
\end{axis}
\end{tikzpicture}
\caption{}
\end{subfigure}\hfill
\begin{subfigure}[b]{0.49\textwidth}\centering
\begin{tikzpicture}
\begin{axis}[graf,width=6.9cm,height=5.4cm,xmin=0.4,xmax=3.2,ymin=0,ymax=1.9,
    xtick={1,2.4},xticklabels={$S_1$,$S_2$},ytick={1,1.5},yticklabels={$T_c$,$T_h$},
    xlabel={$S$},ylabel={$T$},clip=false]
  \fill[fondomedio] (axis cs:1,1) rectangle (axis cs:2.4,1.5);
  \fill[pattern=north east lines,pattern color=acentoclaro] (axis cs:1,0) rectangle (axis cs:2.4,1);
  \draw[rojo,line width=1.1pt] (axis cs:1,1.5) -- (axis cs:2.4,1.5);
  \draw[acento,line width=1.1pt] (axis cs:2.4,1) -- (axis cs:1,1);
  \draw[navy,line width=1.1pt] (axis cs:2.4,1.5) -- (axis cs:2.4,1);
  \draw[navy,line width=1.1pt] (axis cs:1,1) -- (axis cs:1,1.5);
  \draw[rojo,-{Stealth[length=5pt]}] (axis cs:1.55,1.5) -- (axis cs:1.85,1.5);
  \draw[acento,-{Stealth[length=5pt]}] (axis cs:1.85,1) -- (axis cs:1.55,1);
  \node[font=\scriptsize,anchor=south east,inner sep=2pt] at (axis cs:1,1.5) {1};
  \node[font=\scriptsize,anchor=south west,inner sep=2pt] at (axis cs:2.4,1.5) {2};
  \node[font=\scriptsize,anchor=north west,inner sep=2pt] at (axis cs:2.4,1) {3};
  \node[font=\scriptsize,anchor=north east,inner sep=2pt] at (axis cs:1,1) {4};
  \node[font=\scriptsize,rojo,anchor=south] at (axis cs:1.7,1.55) {$Q_h=T_h\,\Delta S$};
  \node[font=\scriptsize,navy] at (axis cs:1.7,1.25) {$|W|$};
  \node[font=\scriptsize,acento,fill=white,inner sep=1pt] at (axis cs:1.7,0.5) {$|Q_c|=T_c\,\Delta S$};
\end{axis}
\end{tikzpicture}
\caption{}
\end{subfigure}
\caption{Ciclo de Carnot de un gas perfecto. (a) Plano del trabajo: dos isotermas (roja, caliente; azul, fría)
y dos adiabáticas reversibles; el área encerrada es $|W|$. (b) Plano del calor: las adiabáticas son verticales
y el ciclo es un rectángulo; el área bajo cada isoterma es su calor.}
\label{fig:carnot}
\end{figure}

\begin{observacion}[Temperatura termodinámica]
Todos los motores reversibles entre los mismos focos tienen el mismo rendimiento, así que
$(Q_c/Q_h)_{\mathrm{rev}}=-T_c/T_h$ no depende de la sustancia. Esto define la escala absoluta de temperaturas
sin referirse a ningún termómetro: $T_c=T_h\,|Q_c/Q_h|_{\mathrm{rev}}$.
\end{observacion}

\begin{ejemplo}[Límites entre \SI{600}{K} y \SI{300}{K}]
Un motor tiene $\eta\leq1-\frac{300}{600}=0{,}50$: con $Q_h=\SI{1}{kJ}$ entrega como mucho \SI{0,50}{kJ}. Un
refrigerador, $\varepsilon\leq\frac{300}{300}=1$; una bomba de calor, $\nu\leq\frac{600}{300}=2$. La igualdad exige
que la máquina y sus intercambios con los focos sean reversibles.
\end{ejemplo}

% =====================================================================
\section{Ecuaciones prácticas}
% =====================================================================

Uniendo la primera ley con $\dbar Q_{\mathrm{rev}}=T\,\dd S$ se obtiene la \emph{relación fundamental}
\begin{equation}\label{eq:fundamental}
  \dd U=T\,\dd S-p\,\dd V,\qquad T=\pd USV,\qquad p=-\pd UVS .
\end{equation}
Es una identidad entre estados de equilibrio, así que vale aunque el proceso real sea irreversible.

\begin{proposicion}[Ecuaciones $T\,\dd S$]
Para un sistema hidrostático,
\begin{equation}\label{eq:TdS}
  \dd S=\frac{C_V}{T}\,\dd T+\pd pTV\dd V,\qquad
  \dd S=\frac{C_p}{T}\,\dd T-\pd VTp\dd p,
\end{equation}
es decir, los coeficientes calorimétricos del Tema 2 valen $l=T\pdi pTV$ y $h=-T\pdi VTp$, y
\begin{equation}\label{eq:CpCVgeneral}
  C_p-C_V=T\pd pTV\pd VTp=\frac{TV\alpha^2}{\kappa_T}\geq0 .
\end{equation}
\end{proposicion}

\begin{demostracion}
La primera es $\dd S=\dbar Q_{\mathrm{rev}}/T=(C_V\dd T+l\,\dd V)/T$ con $l=\pdi UVT+p=T\pdi pTV$ por
\eqref{eq:energetica}. La segunda sale igual con $\dbar Q_{\mathrm{rev}}=C_p\dd T+h\,\dd p$ e imponiendo la
exactitud de $\dd H=T\,\dd S+V\dd p$, que da $h=-T\pdi VTp$. Por último, $C_p-C_V=l\pdi VTp$ (Tema 2), con
$\pdi pTV=\alpha/\kappa_T$ y $\pdi VTp=\alpha V$.
\end{demostracion}

\begin{observacion}[Qué aportan]
Permiten calcular $\Delta S$ y la adiabática ($\dd S=0$) de \emph{cualquier} sistema con solo su ecuación de
estado y $C_V$ (problema~39), y dan la forma general de $C_p-C_V$, siempre no negativa: el agua a \SI{4}{\celsius},
con $\alpha=0$, tiene $C_p=C_V$.
\end{observacion}

% =====================================================================
\section{Ciclos Otto y Joule--Brayton}
% =====================================================================

Ambos ciclos usan un gas perfecto (aire estándar) con dos adiabáticas reversibles; difieren en los tramos con
calor: isócoros en Otto (motor de explosión) e isóbaros en Brayton (turbina de gas). Como $Q_{12}=Q_{34}=0$,
basta con los dos calores:
\begin{equation}\label{eq:otto}
\begin{aligned}
  &\text{Otto:} && Q_{\mathrm{in}}=C_V(T_3-T_2), && Q_{\mathrm{out}}=C_V(T_1-T_4),\\
  &\text{Brayton:} && Q_{\mathrm{in}}=C_p(T_3-T_2), && Q_{\mathrm{out}}=C_p(T_1-T_4).
\end{aligned}
\end{equation}
En los dos, las adiabáticas dan $T_2/T_1=T_3/T_4$ (o, equivalentemente, $\Delta S_{\mathrm{ciclo}}=0$ da
$T_3T_1=T_2T_4$), y
\begin{equation}\label{eq:etaotto}
  \eta=1-\frac{T_4-T_1}{T_3-T_2}=1-\frac{T_1}{T_2}
  =\begin{cases}1-r^{1-\gamma}, & r=V_1/V_2\ \text{(Otto)},\\[2pt] 1-r_p^{(1-\gamma)/\gamma}, & r_p=p_2/p_1\ \text{(Brayton)}.\end{cases}
\end{equation}
Para $r=8$ y $\gamma=1{,}4$, $\eta_{\mathrm{Otto}}=1-8^{-0,4}=0{,}565$.

\begin{figure}[htbp]
\centering
\begin{subfigure}[b]{0.49\textwidth}\centering
\begin{tikzpicture}
% Otto: V1=4, V2=1 (r=4), gamma=1.4 ; p1=1 -> p2=4^1.4=6.964 ; p3=12, p4=12*4^-1.4=1.723
\begin{axis}[graf,width=6.9cm,height=5.2cm,xmin=0.4,xmax=4.8,ymin=0,ymax=13.5,
    xtick={1,4},xticklabels={$V_2$,$V_1$},ytick=\empty,xlabel={$V$},ylabel={$p$}]
  \addplot[navy,line width=1pt,domain=1:4,samples=60] {(4/x)^1.4};
  \addplot[navy,line width=1pt,domain=1:4,samples=60] {12*(1/x)^1.4};
  \addplot[rojo,line width=1.2pt] coordinates {(1,6.964) (1,12)};
  \addplot[acento,line width=1.2pt] coordinates {(4,1.723) (4,1)};
  \fill[black] (axis cs:4,1) circle (1.5pt) node[anchor=north west,font=\scriptsize,inner sep=2.5pt] {1};
  \fill[black] (axis cs:1,6.964) circle (1.5pt) node[anchor=east,font=\scriptsize,inner sep=2.5pt] {2};
  \fill[black] (axis cs:1,12) circle (1.5pt) node[anchor=south west,font=\scriptsize,inner sep=2.5pt] {3};
  \fill[black] (axis cs:4,1.723) circle (1.5pt) node[anchor=west,font=\scriptsize,inner sep=2.5pt] {4};
  \node[font=\scriptsize,rojo,anchor=east] at (axis cs:0.95,9.5) {$Q_{\mathrm{in}}$};
  \node[font=\scriptsize,acento,anchor=west] at (axis cs:4.05,3.0) {$Q_{\mathrm{out}}$};
\end{axis}
\end{tikzpicture}
\caption{}
\end{subfigure}\hfill
\begin{subfigure}[b]{0.49\textwidth}\centering
\begin{tikzpicture}
% Brayton: p1=1, p2=6 ; V1=4 -> V2=4*6^(-1/1.4)=1.111 ; V3=2.2 ; V4=2.2*6^(1/1.4)=7.92
\begin{axis}[graf,width=6.9cm,height=5.2cm,xmin=0,xmax=9,ymin=0,ymax=7.5,
    xtick=\empty,ytick={1,6},yticklabels={$p_1$,$p_2$},xlabel={$V$},ylabel={$p$}]
  \addplot[navy,line width=1pt,domain=1.111:4,samples=60] {(4/x)^1.4};
  \addplot[navy,line width=1pt,domain=2.2:7.92,samples=60] {6*(2.2/x)^1.4};
  \addplot[rojo,line width=1.2pt] coordinates {(1.111,6) (2.2,6)};
  \addplot[acento,line width=1.2pt] coordinates {(7.92,1) (4,1)};
  \fill[black] (axis cs:4,1) circle (1.5pt) node[anchor=north,font=\scriptsize,inner sep=2.5pt] {1};
  \fill[black] (axis cs:1.111,6) circle (1.5pt) node[anchor=south east,font=\scriptsize,inner sep=2.5pt] {2};
  \fill[black] (axis cs:2.2,6) circle (1.5pt) node[anchor=south west,font=\scriptsize,inner sep=2.5pt] {3};
  \fill[black] (axis cs:7.92,1) circle (1.5pt) node[anchor=north,font=\scriptsize,inner sep=2.5pt] {4};
  \node[font=\scriptsize,rojo,anchor=south] at (axis cs:1.65,6.45) {$Q_{\mathrm{in}}$};
  \node[font=\scriptsize,acento,anchor=south] at (axis cs:6.0,1.15) {$Q_{\mathrm{out}}$};
\end{axis}
\end{tikzpicture}
\caption{}
\end{subfigure}
\caption{(a) Ciclo Otto: el calor entra y sale a volumen constante. (b) Ciclo Joule--Brayton: entra y sale a
presión constante. En ambos, $1\to2$ y $3\to4$ son adiabáticas reversibles.}
\label{fig:ottobrayton}
\end{figure}

% =====================================================================
\section{Trabajo perdido y trabajo máximo}
% =====================================================================

\begin{proposicion}[Trabajo perdido]
Para un motor entre dos focos, con $S_{\mathrm{gen}}=-Q_h/T_h-Q_c/T_c\geq0$,
\begin{equation}\label{eq:perdido}
  |W|=\Bigl(1-\frac{T_c}{T_h}\Bigr)Q_h-T_cS_{\mathrm{gen}},\qquad
  W_{\mathrm{perdido}}\equiv|W|_{\max}-|W|=T_cS_{\mathrm{gen}} .
\end{equation}
\end{proposicion}

\begin{demostracion}
Despejando $Q_c=-T_c(S_{\mathrm{gen}}+Q_h/T_h)$ y llevándolo a $|W|=Q_h+Q_c$ se obtiene la primera expresión;
el máximo, con $S_{\mathrm{gen}}=0$, es el de Carnot.
\end{demostracion}

\begin{ejemplo}[Cuánto trabajo se pierde]
Un motor entre 600 y \SI{300}{K} absorbe \SI{1}{kJ} y entrega \SI{0,40}{kJ}, luego cede \SI{0,60}{kJ}:
$S_{\mathrm{gen}}=-\frac{1000}{600}+\frac{600}{300}=\SI{0,333}{J/K}$ y $W_{\mathrm{perdido}}=300\cdot\tfrac13=\SI{100}{J}$, la
diferencia con los \SI{0,50}{kJ} de Carnot.
\end{ejemplo}

\begin{observacion}[Trabajo máximo con cuerpos finitos]
Si los focos son cuerpos de capacidad finita, su temperatura cambia y no hay un $\eta_C$ fijo. El trabajo máximo
se obtiene igualmente imponiendo reversibilidad, $\Delta S_{\mathrm{univ}}=0$, que fija la temperatura final; la
energía sale después de la primera ley (problema~38).
\end{observacion}
